\section{Methodology}\label{sec:method}

We frame berth scheduling as a decision made under an uncertain service time and
ask how a predictor of that service time should be trained. The scheduler is a
discrete berth allocation model (Section~\ref{sec:milp}) whose only uncertain
input is the berth-occupation time of each vessel. A predictor supplies that
input from pre-berthing features (Section~\ref{sec:predictors}). The
methodological contribution is the training objective for that predictor: rather
than fit the service time in isolation, we train it against the downstream cost
of the schedule it induces (Section~\ref{sec:dfl}). The mechanism that motivates
this choice is a \emph{cascade asymmetry}. Under-predicting a vessel's service
time leads the optimizer to pack the following vessel too early; when the true
service time is realized, the overlap propagates as a chain of delays down the
berth queue. Over-predicting only leaves a berth briefly idle and does not
cascade. A squared-error loss treats the two errors symmetrically. A
decision-focused loss does not, because it scores the realized schedule.

\subsection{Discrete berth allocation model}\label{sec:milp}

We adopt the discrete formulation of the berth allocation problem (BAP), in which
the quay is partitioned into a finite set of berths and at most one vessel
occupies a berth at any time \citep{imai2001dynamic, buhrkal2011models}. The
discrete variant matches the terminal under study, whose quay is operated as a
fixed set of numbered berths with vessel-specific compatibility, and it keeps the
assignment and precedence structure that the decision-focused gradients later
exploit. Let
$I=\{1,\dots,N\}$ be the vessels to be scheduled in the planning week and
$B=\{1,\dots,B\}$ the berths. Each vessel $i$ is compatible with a subset
$B_i\subseteq B$ of berths, determined by draft, length, and cargo type; an
assignment to $b\notin B_i$ is forbidden, which we impose by restricting the
assignment variables rather than through side constraints
\citep{alzaabi2016berth}. A distinguished subset $S\subseteq I$ of priority
service vessels carries fixed reservations and is excluded from the optimized
objective. Each vessel $i$ has an arrival time $a_i$ and a berth-occupation
(service) time $\tau_i$. For $i\notin S$ the service time is \emph{predicted},
$\tau_i=\hat\tau_i$; for a reserved vessel $r\in S$ it is fixed at its declared
duration $D_r$. A single navigation channel of transit time $c$ may connect the
anchorage to the berths; setting $c=0$ recovers the berth-only model.

The decision variables are the assignment $x_{ib}\in\{0,1\}$, defined only for
$b\in B_i$; the berth start time $s_i\ge a_i$; and the pairwise precedence
$z_{ijb}\in\{0,1\}$, equal to $1$ when vessel $i$ precedes vessel $j$ at berth $b$.
When the channel is active we add channel-entry and channel-exit times
$e^{\text{in}}_i,e^{\text{out}}_i$ and a transit-order variable $y_{ij}$.

The default objective minimizes total unweighted waiting time,
\begin{equation}\label{eq:obj}
\min \; \sum_{i\in I\setminus S}\bigl(s_i-a_i\bigr),
\end{equation}
where the channel-entry time $e^{\text{in}}_i$ replaces $s_i$ when the channel is
active. An equivalent berth-idle form, $\min\sum_{b\in B}(\mathrm{LC}_b-\mathrm{FS}_b)$
over the last completion $\mathrm{LC}_b$ and first start $\mathrm{FS}_b$ at each
berth, yields the same optimizer. Because arrival times are data, unweighted
waiting equals total completion $\sum_{i}(s_i+\tau_i)$ up to an additive
constant; we exploit this identity in Section~\ref{sec:dfl}. The objective is
deliberately unweighted. Priority weights, which would place a per-vessel cost
vector in front of the completion terms, are left to future work.

The constraints are as follows. Each vessel is assigned to exactly one compatible
berth,
\begin{equation}\label{eq:assign}
\sum_{b\in B_i} x_{ib}=1 \qquad \forall i\in I .
\end{equation}
A vessel cannot start before it arrives,
\begin{equation}\label{eq:release}
s_i \ge a_i \qquad \forall i\in I .
\end{equation}
No two vessels overlap at a shared berth. For each pair $\{i,j\}$ and each berth
$b\in B_i\cap B_j$,
\begin{align}
z_{ijb}+z_{jib} &\le 1, \label{eq:prec-one}\\
z_{ijb}+z_{jib} &\ge x_{ib}+x_{jb}-1, \label{eq:prec-active}\\
s_j &\ge s_i + \tau_i - M\,(1-z_{ijb}), \label{eq:prec-bigm}
\end{align}
so that whenever $i$ and $j$ share berth $b$ exactly one ordering is selected
(\ref{eq:prec-one})--(\ref{eq:prec-active}) and the successor cannot start until
the predecessor's service is complete (\ref{eq:prec-bigm}). Priority
reservations hold a fixed berth interval with no backfill,
\begin{equation}\label{eq:reserve}
s_r = R_r, \quad \tau_r = D_r \qquad \forall r\in S .
\end{equation}
The optional shared channel enforces the transit time and prevents two transits
from overlapping,
\begin{align}
s_i &\ge e^{\text{in}}_i + c, \qquad e^{\text{out}}_i \ge s_i + \tau_i, \label{eq:chan-time}\\
e^{\text{in}}_j &\ge e^{\text{out}}_i - M\,(1-y_{ij}), \qquad y_{ij}+y_{ji}\le 1, \label{eq:chan-order}
\end{align}
which reduces to the berth-only model at $c=0$.

Constraints (\ref{eq:prec-bigm}) and (\ref{eq:chan-time})--(\ref{eq:chan-order})
carry the methodological weight of this section. The predicted service time
$\hat\tau_i$ does not enter the model as a linear cost coefficient in the
objective; it enters the big-$M$ \emph{constraints} that fix when the following
vessel may start. A prediction error therefore shifts the feasible region rather
than merely rescaling the objective, which is precisely the setting that dictates
the gradient choice in Section~\ref{sec:dfl}. The precedence and completion
structure follows the discrete BAP family of \citet{cordeau2005berth} and
\citet{buhrkal2011models}. The big-$M$ is sized as $\max_i a_i + N\cdot\bar\tau$,
where $\bar\tau$ is an upper bound on any single service time, so that
$M$ dominates the worst-case completion of the queue and never truncates a
feasible schedule. The model is built in Pyomo and solved with Gurobi to proven
optimality (MIPGap set to $0$), so that any suboptimality in the realized
schedule is attributable to the predicted input rather than to the solver.

\subsection{Service-time prediction}\label{sec:predictors}

The predictor $f_\theta$ maps a vector of pre-berthing features $x$, available
before the vessel is scheduled, to an estimated service time $\hat\tau=f_\theta(x)$.
We benchmark a model zoo under $5$-fold cross-validation stratified by berth
(\emph{Sitio}), so that each fold preserves the berth composition of the full
sample. Hyperparameters are tuned with Optuna using the tree-structured Parzen
estimator, allotting $100$ trials to the tree ensembles and $40$ to the neural
models. A fresh preprocessor is fit on the training partition of each fold and
applied to the held-out partition, so no target or scaling information leaks
across the split.

The zoo spans four families. A linear ridge-regularized head serves both as a
baseline and as the shared backbone for the decision-focused experiments below.
The gradient-boosted and bagged trees comprise random forest, XGBoost, and
LightGBM. The tabular deep-learning models are RealMLP \citep{holzmuller2024realmlp},
TabM \citep{gorishniy2025tabm}, and NODE \citep{popov2020node}. Global-mean and
group-mean (per-berth) predictors provide the naive reference points. For the
controlled comparison between predict-then-optimize and decision-focused learning
in Section~\ref{sec:dfl}, only the linear head is carried into both arms. Holding
the architecture fixed across arms ensures that any difference in downstream
schedule cost is attributable to the training loss and not to model capacity.

\subsection{Decision-focused learning}\label{sec:dfl}

The conventional pipeline is predict-then-optimize (PtO): fit $f_\theta$ by
minimizing the mean squared error
\begin{equation}\label{eq:mse}
\mathcal{L}_{\text{PtO}}(\theta) = \sum_{i\in I}\bigl\|\hat\tau_i-\tau_i\bigr\|^2 ,
\end{equation}
then pass the point predictions to the MILP of Section~\ref{sec:milp}. The two
stages are trained and executed independently, so the predictor never observes
the schedule its errors produce.

We measure schedule quality by \emph{regret}. Given predictions $\hat\tau$, we
solve the MILP to obtain a decision $(x,z)$, that is, the berth assignment and
the precedence order. We then lock that assignment and order and re-derive
feasible start times under the \emph{true} service times $\tau$ by greedy
earliest-start packing along each berth's fixed order, and score the realized
schedule at its true completion. This evaluation mirrors terminal practice: once
a vessel is committed to a berth and a position in the queue, the assignment and
order are held while the actual start times float to absorb the realized service
durations. The realized cost is

\begin{equation}\label{eq:cost}
C = \sum_{i\in I}\bigl(s_i+\tau_i\bigr).
\end{equation}
Regret is the excess of this realized cost over the full-information cost,
\begin{equation}\label{eq:regret}
\mathrm{Regret} = C_{\text{method}} - C_{\text{FI}},
\end{equation}
where $C_{\text{FI}}$ is obtained by solving the MILP directly under the true
$\tau$. Regret is nonnegative under the present objective. Because unweighted
waiting equals total completion (\ref{eq:cost}) up to a constant, the
full-information decision is the exact minimizer of the scored cost, so every
alternative decision re-derives to a cost of at least $C_{\text{FI}}$. We
verified this property both analytically and by brute-force enumeration over
small instances. We note that the nonnegativity of regret is tied to the
unweighted objective and would require re-examination once priority weights are
introduced.

Decision-focused learning (DFL) replaces the surrogate loss (\ref{eq:mse}) with
the realized schedule cost itself. Let $s^\star$ denote the start times the solver
produces when it optimizes under the predicted $\hat\tau$. The DFL loss evaluates
those solver-chosen start times at the true service times,
\begin{equation}\label{eq:dfl}
\mathcal{L}_{\text{DFL}}(\theta) = \sum_{i\in I}\bigl(s^\star_i+\tau_i\bigr),
\end{equation}
and is backpropagated through the MILP into $\theta$. The loss rewards
predictions that yield schedules which remain cheap once reality is realized,
which is exactly the property the cascade asymmetry makes MSE blind to.

The gradient of (\ref{eq:dfl}) must pass through the argmin of a mixed-integer
program. Because $\hat\tau$ enters the big-$M$ \emph{constraints}
(\ref{eq:prec-bigm}) and (\ref{eq:chan-time})--(\ref{eq:chan-order}) rather than
a linear objective cost vector, the SPO+ surrogate \citep{elmachtoub2022spo},
which is derived for uncertain objective coefficients, does not apply directly.
We therefore use solver-differentiation methods that require only black-box
solver calls on perturbed inputs. Blackbox differentiation
\citep{pogancic2020blackbox} constructs an informative gradient from a single
interpolation strength $\lambda$ and two solves per training step. Perturbed
optimizers \citep{berthet2020learning} smooth the argmin with Gaussian
perturbations of the input and average over $K$ samples. Both are implemented
through PyEPO \citep{tang2024pyepo}. The alternative family for
predictions-in-constraints is implicit differentiation through relaxed
optimization layers \citep{amos2017optnet, wilder2019melding, ferber2020mipaal};
we did not adopt it here on grounds of computational cost and scalability to the
discrete BAP. For the broader landscape of decision-focused training we refer to
the surveys of \citet{mandi2024decision} and \citet{kotary2021end}, and for the
specific and less-developed case of predicted parameters entering constraints to
\citet{hu2023twostage} and \citet{mandi2025feasibility}. Each DFL model is
warm-started from the trained PtO weights of the shared linear head, so the only
difference between the two arms is the training loss.
